% Generated draft: requires mathematical and source review.
\documentclass[11pt]{article}
\usepackage[margin=25mm]{geometry}
\usepackage{fontspec}
\setmainfont{Latin Modern Roman}
\usepackage{amsmath,amssymb,mathtools,unicode-math}
\setmathfont{Latin Modern Math}
\usepackage{enumitem,needspace,xurl,hyperref}
\setlength{\parindent}{0pt}
\setlength{\parskip}{5pt}
\setlength{\emergencystretch}{3em}
\newcommand{\N}{\mathbb{N}}
\newcommand{\Z}{\mathbb{Z}}
\newcommand{\Q}{\mathbb{Q}}
\newcommand{\R}{\mathbb{R}}
\newcommand{\C}{\mathbb{C}}
\newcommand{\F}{\mathbb{F}}
\newcommand{\A}{\mathbb{A}}
\newcommand{\PP}{\mathbb{P}}
\newcommand{\E}{\mathbb{E}}
\newcommand{\eps}{\varepsilon}
\newcommand{\dd}{\,\mathrm{d}}
\begin{document}
\section*{UnsolvedMath ID: 10400174}
\textbf{Source problem code:} catalogue entry
\subsection{Definitions and mathematical statement}
A subfactor is an inclusion of von Neumann factors, with standard-invariant/bimodule data from which suitable \(6j\)-symbols may be obtained. N. Sato asks to construct a well-defined state-sum-type invariant from a \emph{strongly amenable} subfactor. The technical strong-amenability definition, index/type hypotheses, label data and coefficient conventions are not supplied.

In the intended state-sum programme, a triangulated manifold is evaluated by summing weights associated to admissible labelings, built from the subfactor data. A well-defined construction must specify the weights and normalization and justify the relevant triangulation invariance; when the labels or sums are infinite, convergence or a specified regularization also requires justification. The source does not restrict to finite-depth subfactors, so replacing strong amenability by finite depth would narrow the task.

The source explicitly contrasts the problem with Ponzano--Regge: it reports no general asymptotic description of the quantum \(6j\)-symbols here, while recalling an asymptotic description for \(\operatorname{SU}(2)\) \(6j\)-symbols. This is a historical qualification, not a new nonexistence theorem for asymptotic methods.

The appended context considers the Turaev--Viro--Ocneanu invariant of a closed 3-manifold from a subfactor and reports that the Longo--Rehren construction relates it to a Reshetikhin--Turaev-type invariant. When the starting data have nondegenerate braiding, it reports splitting into a Reshetikhin--Turaev invariant and its complex conjugate, conventionally a product \(Z_{\mathrm{RT}}(M)\overline{Z_{\mathrm{RT}}(M)}\), with normalization hypotheses still to be supplied. This contextual statement must not be read as solving the construction for every strongly amenable subfactor.

Finally, the excerpt says the following question would aid a theory of minimal nondegenerate extensions of degenerate braidings, but does not state that next question. It is retained as motivation rather than appended as an invented additional theorem.
\subsection{Short English statement}
Construct a well-defined state sum from strongly amenable subfactors, retaining possible infinite-label and asymptotic difficulties. Preserve the source's Longo--Rehren and braided factorization context without treating it as a full solution.
\subsection{Sources}
\begin{itemize}
\item Local UnsolvedMath registry, permanent ID 10400174, source code catalogue entry; statement attribution N. Sato; proposer field and proposal year absent; update 2026-08-06.
\item Ohtsuki, \emph{Problems on invariants of knots and 3-manifolds} (2002), Problem 9.10, PDF page 141: \url{https://msp.org/gtm/2002/04/gtm-2002-04-024s.pdf}. Registry access date 2026-07-29; PDF font-structure extraction. Provenance distinguishes 231 numbered bold headings. Present primary-literature status and redistribution terms require review.
\item The source names Ponzano--Regge, Turaev--Viro--Ocneanu, Longo--Rehren and Reshetikhin--Turaev constructions without exact bibliography entries in this excerpt. The research summary gives no verified resolving reference.
\end{itemize}
\end{document}
